\documentclass[11pt]{article}

\usepackage[a4paper,margin=1in]{geometry}
\usepackage{amsmath,amssymb,booktabs,graphicx}
\usepackage[numbers,sort&compress]{natbib}
\usepackage[colorlinks=true,linkcolor=blue,citecolor=blue,urlcolor=blue]{hyperref}

% Wrapper for generated result macros. Build with `make paper-assets` before compiling.
\IfFileExists{../reports/paper/latex/result_macros.tex}{%
  \input{../reports/paper/latex/result_macros.tex}%
}{%
  \input{../../reports/paper/latex/result_macros.tex}%
}


\title{From Extreme Samples to Failure Episodes: A Reproducible Dynamical-EVT Diagnostic Workflow for Predictive Maintenance}
\author{Diogo Ribeiro}
\date{\today}

\begin{document}
\maketitle

\begin{abstract}
Industrial predictive-maintenance telemetry is commonly evaluated through pointwise anomaly scores, although maintenance decisions concern failure episodes, repair intervals, and alarm burden. This paper reports a reproducible workflow that converts causal machine-state observations into regime-conditioned extreme-event diagnostics, alarm episodes, and artifact-backed claims. The implementation combines dangerous-region observables, threshold and run-length sensitivity analysis, extremal-index diagnostics, event-level matching, baseline provenance, and a frozen evaluation protocol. The completed asset bundle records \DynEvtClaimCount{} claim-ledger entries, of which \DynEvtSupportedClaimCount{} are permitted for manuscript use under the generated ledger. The present evidence supports a narrow conclusion: the workflow can expose estimator sensitivity, clustering, calibration diagnostics, and operational alarm trade-offs from generated artifacts, while it does not support a population-level superiority claim for industrial predictive maintenance. All reported claims are linked to experiment identifier \DynEvtExperimentId{}, protocol hash \DynEvtProtocolHash{}, and generated figures or tables.
\end{abstract}

\section{Introduction}
\label{sec:introduction}

Industrial maintenance alarms are useful only when they map to the decision unit that operators and engineers actually face. In some settings the unit is a failure episode, in others a repair history, a laboratory component-condition cycle, or a production example with a downstream yield label. Standard pointwise anomaly detection can be accurate at the sample level while still producing a poor operational object. A burst of extreme samples near one physical failure may be counted as many alarms; a vehicle-level repair problem may be misread as timestamp-level warning; and millions of telemetry rows can obscure the small number of independent events or entities that actually support inference.

Cyclic machinery creates an additional statistical problem. Compressor, engine, and actuator signals are shaped by operating regimes, control loops, maintenance states, and sensor preprocessing. Extreme observations therefore need not be independent exceedances of a stationary distribution. Dynamical extreme-value theory connects high observations of functions evaluated along a trajectory with returns to small target sets, and it explains why periodic or recurrent structure can create clustered extremes \citep{Freitas2008Hitting,Freitas2012Periodicity}. Under finite noisy observations, the same asymptotic quantities may be masked or distorted \citep{Faranda2013Noise}. These facts make naive threshold exceedance counts an unstable basis for maintenance claims.

This paper studies a constrained contribution: a reproducible diagnostic workflow that turns machine-state trajectories and tabular condition-monitoring examples into artifact-backed evidence, with explicit warnings about what the artifacts do not prove. The workflow defines causal state vectors, operating regimes, training-defined dangerous regions, regime-conditioned thresholds, exceedance clusters, lagged multivariate patterns, alarm episodes, and dataset-specific condition diagnostics evaluated under a frozen protocol. The central question is not whether a new model is a general winner across industrial settings. The question is whether generated artifacts make dependence, clustering, alarm burden, sensitivity, calibration, estimand changes, and unsupported claims visible enough to constrain a defensible predictive-maintenance paper.

The contributions are therefore deliberately bounded. First, the implementation provides a dynamical-EVT-inspired formulation for regime-conditioned extreme diagnostics in predictive maintenance. Second, it produces finite-sample simulation and sensitivity artifacts for extremal-index recovery and threshold/run-length choices. Third, it runs a five-dataset real-data matrix covering compressor telemetry, fleet repair histories, hydraulic component-condition cycles, and semiconductor process/yield examples. Fourth, it freezes evaluation and audit protocols so recall, false alarms, duplicate alarms, lead time, calibration, verification status, and industrial diagnostic rows are generated rather than hand-curated. Finally, it generates a claim ledger that allows only artifact-backed statements in the manuscript and provides build commands that regenerate the figures, tables, provenance, manuscript, supplement, and audits from version-controlled code. These contributions correspond to claim-ledger entries CLM-001 through CLM-007; claims outside that ledger are not used as positive evidence.

\section{Related Work}
\label{sec:related}

Predictive-maintenance models often begin with supervised classification, reconstruction error, distance-based anomaly detection, or streaming thresholds. These methods are useful engineering baselines, but their common output is a timestamped score rather than an episode-level maintenance decision. EVT-based streaming detectors such as SPOT use peaks-over-threshold modelling to adapt anomaly thresholds in streams \citep{Siffer2017SPOT}; they address threshold calibration, but they do not by themselves resolve operating-regime mixture, duplicate alarms per failure, or the distinction between clustered observations and independent maintenance events.

Classical declustering and extremal-index estimation address dependence among extremes. Ferro and Segers proposed an inference scheme for clusters of extreme values using inter-exceedance times and the extremal index \citep{FerroSegers2003Clusters}. In dynamical systems, hitting-time statistics connect returns to target neighborhoods with extreme-value laws \citep{Freitas2008Hitting}. Periodic points can generate an extremal index smaller than one, giving a mathematical route from recurrence to clustered extremes \citep{Freitas2012Periodicity}. The finite-sample industrial setting is less clean: noise may make deterministic recurrence appear closer to independence, and observed clustering can also arise from smoothing, missingness, persistence, or regime mixture \citep{Faranda2013Noise}. Recent multivariate work extends dynamical EVT to cross-sectional dependence among vector-valued observables \citep{Aimino2024Multivariate}, which motivates the lag-pattern diagnostics used here.

The industrial datasets used by the repository are public predictive-maintenance resources with different independence units. MetroPT is an air-production-unit telemetry dataset for predictive maintenance \citep{Veloso2022MetroPT}. MetroPT2 is a replication-style benchmark released through Zenodo \citep{MetroPT2Zenodo2023}. SCANIA Component X is a fleet-level multivariate time-series dataset with operational histories and repair records for an anonymized component \citep{Kharazian2025Scania}. These datasets support useful case studies, but they do not remove the need for event-level accounting: one long fault interval is not equivalent to thousands of independent statistical units. The methodological gap addressed here is the absence of a reproducible layer that connects dynamical extreme diagnostics to alarm episodes, sensitivity analyses, and claim constraints.

\section{Method}
\label{sec:method}

Let \(Z_t\in\mathbb{R}^p\) be the causal state vector at sample \(t\). For compressor telemetry, \(Z_t\) contains the declared raw pressure, temperature, current, and flow channels used by the real-data diagnostic. For tabular condition-monitoring adapters, \(Z_t\) contains the declared feature vector for the cycle, vehicle, or production example.

Scaling is fitted on the training partition only. Each feature is centred by its training median and divided by a scale
\[
    s_j=\max\!\left\{\mathrm{IQR}_j,\ \kappa\left(q_{0.999,j}-q_{0.001,j}\right)\right\},\qquad \kappa=0.01,
\]
where all quantiles are training quantiles. The floor is not cosmetic. An interquartile range measures the spread of the central half, and a channel that idles at a fixed value between excursions has a central half made of quantisation noise: MetroPT's \texttt{pressure\_tp2} has an interquartile range of \(0.004\) against a \(0.1\)-to-\(99.9\) percentile range near ten, so its central half occupies four hundredths of one percent of its range. Dividing by that spread multiplies ordinary variation by roughly two hundred and fifty, and in a maximum-absolute-\(z\) aggregate one such channel then sets the score, its threshold and its exceedances for every other channel. Flooring the denominator at a fraction of the feature's own range leaves well-behaved features untouched, since for them the interquartile range already exceeds \(\kappa\) times the range, and prevents a quantised or mostly idle channel from setting the scale for the rest. The range is measured between the \(0.1\)st and \(99.9\)th percentiles rather than the 1st and 99th because a channel whose excursions are rarer than one percent has a narrow 1-to-99 range as well, which would leave the floor as degenerate as the quantity it bounds.

Operating regimes \(R_t\), when used in the sensitivity assets, are inferred without test labels. Regimes are conditioning variables for threshold estimation and are not interpreted as mechanical causes. The registered real-data matrix reported in Table~\ref{tab:event-benchmark-compact} and Supplementary Table~\ref{X-tab:non-event-results} uses a conservative train-split robust score rather than the full dangerous-region construction.

The implementation also provides dangerous-region scoring for separate diagnostics. The failure-proximal prototype variant uses only training failure intervals and records the source interval for every prototype. The unsupervised rare-state variant ignores failure labels and fits a high-distance rare-state set from training data only. These variants are available for diagnostic analysis and ablation coverage, but the registered real-data matrix in this manuscript is the conservative high-quantile score; it should not be read as evidence that both dangerous-region variants have already been optimized on every dataset.

For a target region \(A\), the scalar observable is
\[
    X_t=-\log\{d(Z_t,A)+\epsilon\},
\]
where \(d\) is the robust Euclidean distance in scaled coordinates and \(\epsilon>0\) prevents singularities. Larger \(X_t\) means closer proximity to the training-defined region. This is an anomaly or recurrence score, not a calibrated probability of failure.

For each regime \(r\), a threshold \(u_r(q)\) is estimated from the declared training or calibration partition at quantile \(q\). Exceedances are \(I_t(q,r)=\mathbf{1}\{X_t>u_r(q),R_t=r\}\). Runs declustering merges exceedances separated by at most \(m\) samples. The runs estimator, Ferro--Segers inter-exceedance estimator, and K-gaps-style declustering diagnostic summarize clustering under the frozen \(q\) and \(m\) settings \citep{FerroSegers2003Clusters,SuvegesDavison2010KGap}. The inverse extremal index is discussed only as an asymptotic cluster-size diagnostic under the assumptions of the estimator.

Pointwise exceedances are converted into alarm episodes by the declared merge gap. A test alarm is a true event alarm only if the episode intersects the predeclared warning window for a held-out failure under the frozen matching tolerance. All event recall, event precision, false-alarm burden, duplicate-alarm count, lead time, time under warning, and utility metrics are computed after this episode conversion when event labels exist. This prevents timestamp counts from being treated as independent maintenance evidence.

\paragraph{What is being evaluated.} Table~\ref{tab:method-taxonomy} states, for each of the nineteen benchmarked methods, what its score measures, how its threshold is chosen, how exceedances become episodes, which partition its parameters are fitted on, and where extreme-value theory enters. Three facts in that table matter for how this paper should be read, and none of them is apparent from the method names.

First, the recurrence observable \(X_t=-\log\{d(Z_t,A)+\epsilon\}\) defined above is used by the three target-region methods. It is not the score of the method labelled as registered. That method scores a robust maximum absolute \(z\), which is a multivariate anomaly magnitude, and it shares that score with seven of the baselines; what distinguishes it from them is its event policy.

Second, extreme-value theory changes the alarms a method raises in only five of the nineteen. A tail model sets the threshold for the peaks-over-threshold baseline and for SPOT, and an estimated extremal index sets the declustering run length for the Ferro--Segers policy, the K-gaps policy and the registered method. For the remaining fourteen the extremal-index estimates reported here are diagnostics computed beside the alarms; those methods would raise identical alarms if the estimates were discarded.

Third, and following from the first two, the contribution offered here is the evaluation architecture rather than a detector. What this paper supplies is a layered decomposition that locates where a detector loses its signal, a rank separation and a shift-based null that say whether a match carries information, and an event-level protocol under which simple and elaborate methods are compared on the same terms. The dynamical-EVT material motivates the recurrence observable and supplies the declustering rule, and the empirical results are evidence about what those produce under this protocol on two failures. They are not evidence that a robust anomaly score thresholded and declustered constitutes a dynamical-EVT detector, and the paper does not claim it.

\input{../reports/paper/latex/method_taxonomy.tex}

\paragraph{Warning-window rank separation and its null.} To separate a score that carries no pre-onset ordering from a threshold that discarded the ordering a score did carry, we compute the probability that a sample drawn from the declared horizon before a failure onset outranks one drawn from the rest of the test split. This is the Mann--Whitney statistic on midranks, so it is invariant to any monotone transformation of the score and comparable across methods; \(0.5\) is chance. Midranks matter here because industrial telemetry ties heavily during idle operation, and breaking ties by position would make a constant score appear informative purely because the warning window sits at one end of the series.

A recall of one is only evidence of anticipation if a detector could not have achieved it by alarming often enough. We therefore compute a chance-match probability by circularly shifting the entire alarm series by a uniformly random offset and asking how often at least one episode still intersects a matching window. The shift preserves the episode count, the episode durations, and the burstiness and regime structure of their spacing, and randomises only the alignment to the failure. Because an episode \([s,s+L]\) intersects a window \([a,b]\) exactly when \((s+d)\bmod n\) falls in an interval of length \(b-a+L+1\), the quantity is a measure of a union of circular intervals and is computed exactly rather than sampled.

We do not use a homogeneous Poisson model for this null. Poisson arrivals are independent and occur at a constant rate, whereas these episodes are produced by declustering a temporally dependent score, their rate varies with operating regime, and the merge rule induces further dependence between neighbouring episodes. A null that assumes extremes do not cluster cannot calibrate a paper about clustered extremes.

\paragraph{Event matching.} Let \(P=(p_1,\dots,p_m)\) be the alarm episodes after merging and \(F=(f_1,\dots,f_k)\) the labelled failures. For failure \(f\) with onset \(f^{\mathrm{start}}\), the matching window under horizon \(h\), pre-tolerance \(\tau_-\) and post-tolerance \(\tau_+\) is
\[
    W(f)=\left[\max\{0,\ f^{\mathrm{start}}-h-\tau_-\},\ f^{\mathrm{start}}+\tau_+\right].
\]
The window ends at onset plus \(\tau_+\), not at the end of the failure interval, so an alarm raised well inside a multi-day failure is not counted as a warning. An alarm \(p\) is \emph{eligible} for \(f\) when \(p\cap W(f)\neq\emptyset\); an alarm that begins before the window and extends into it is eligible, since only intersection is required. Matching is one-to-one: each alarm may be assigned to at most one failure and each failure to at most one alarm, chosen by a maximum-cardinality assignment over the eligible pairs, with ties broken toward the longer lead time. Alarm reuse across failures is available in the protocol but disabled here. Every alarm then falls in exactly one of three classes:
\[
\begin{array}{ll}
\text{matched} & \text{assigned to some failure;}\\
\text{duplicate} & \text{eligible for some failure but not assigned;}\\
\text{false} & \text{eligible for no failure.}
\end{array}
\]
Event recall is matched failures over \(|F|\). Event precision is
\[
    \mathrm{precision}=\frac{\#\text{matched}}{\#\text{matched}+\#\text{duplicate}+\#\text{false}}=\frac{\#\text{matched}}{m},
\]
so duplicates \emph{do} enter the denominator and precision is matched alarms over all alarms raised. This is deliberate: an operator responds to a duplicate exactly as to any other alarm, and excluding duplicates would let a detector improve its precision by emitting more alarms near a failure it has already matched.

These conventions determine the reported numbers, and one consequence deserves emphasis. Because episodes are formed by merging before matching, a detector that is permanently on produces a single episode that is matched, with no duplicates and no false alarms, and therefore attains precision and recall of one. Episode-level metrics cannot detect that state, which is why every alarm-burden comparison in this paper reports time under warning alongside episode counts, and why the benchmark flags any method whose warning exposure exceeds half the test period.

\paragraph{Joint utility.} The control comparisons use one scalar, declared before the comparison that consumes it. For a detector with detection indicator \(D\in\{0,1\}\), false alarms per operating day \(a\), warning exposure fraction \(c\), and lead time \(\ell\) in samples,
\[
    U=w_D D-w_a a-w_c c+w_\ell D\,\frac{\min\{\ell,\ell_{\max}\}}{\ell_{\max}},
\]
with \(w_D=100\), \(w_a=1\), \(w_c=10\), \(w_\ell=5\) and \(\ell_{\max}=7200\) samples. The units are therefore: one detection is worth 100, one false alarm per operating day costs 1, being under warning for the whole test period costs 10, and a lead time at or beyond two hours adds 5. The ordering is the substantive claim. Detection dominates by construction, so a method must avoid more than a hundred false alarms per day to justify missing a failure, which no method here approaches; the lead-time reward is capped and paid only on detection, so it cannot make a missed failure look competitive. Section~\ref{sec:results} reports a sensitivity sweep varying each weight by factors from \(0.25\) to \(4\), and the component metrics are reported separately so a reader who prefers different weights can recompute.

The score-stratification target for horizon \(h\) is
\[
    Y_t^{(h)}=\mathbf{1}\{\hbox{a labelled event begins within }h\hbox{ after }t\}.
\]
In the current evidence, model outputs are anomaly ranks or recurrence scores. Probability-calibration language is therefore restricted: Brier-style summaries are reported only as diagnostics of score stratification unless a separate calibration partition contains enough positive independent units for an explicit probability map.

The leakage contract is: training features cannot contain held-out failure samples; scaling, regimes, target regions, thresholds, lag selection, merge-gap selection, and calibration cannot use test events; SCANIA vehicles cannot appear in multiple partitions; and feature windows cannot cross maintenance-reset boundaries. Table~\ref{X-tab:leakage-audit} gives the terminal audit result for each registered check. Under the current registered diagnostic the algorithm is:
\[
\begin{array}{ll}
1. & \text{Split entities or intervals into training, validation, calibration, and test partitions.}\\
2. & \text{Fit scaling and robust score thresholds on allowed partitions only.}\\
3. & \text{Select } q,m,\text{ and merge gap on validation/calibration partitions only where used.}\\
4. & \text{Freeze all parameters and score the held-out test partition causally.}\\
5. & \text{Extract exceedance clusters, merge alarm episodes, and match them to labelled events.}\\
6. & \text{Report event metrics, uncertainty by independent unit, and non-estimable cases.}
\end{array}
\]


The practical complexity for scoring is \(O(npK)\), where \(n\) is the number of samples or entities, \(p\) is the state dimension, and \(K\) is the number of stored prototypes or rare-state summaries. Thresholding and runs declustering are linear in \(n\). The streaming-compatible parts are causal feature construction, scaling with frozen parameters, scoring, thresholding, and alarm merging. Prototype construction, model selection, and final event matching are offline evaluation steps.

\section{Simulation Study}
\label{sec:simulation}

The simulation layer is designed to falsify overconfident interpretations before industrial data are used. The registered systems include independent extremes, periodic dynamics, noisy periodic dynamics, regime mixtures, missingness perturbations, smoothing and downsampling perturbations, and degradation-like cyclic trajectories. Each simulation cell records its configuration, random seed, threshold grid, run-length grid, estimator status, and raw outputs sufficient to regenerate the paper assets.

\subsection{Estimators under comparison}

Six estimators of the extremal index are compared: runs declustering, the Ferro--Segers intervals estimator, K-gaps, the reciprocal mean block-cluster size, the disjoint-blocks estimator, and a no-declustering reference fixed at \(\theta = 1\). The last is not a competitive estimator but a control: any procedure that cannot separate a clustered process from an independent one should be indistinguishable from it. The runs and K-gaps estimators are tuned by the declustering run length; the block-based estimators are tuned by the block size, which is set to \(\lceil \sqrt{n} \rceil\) so that the tuning parameter is declared rather than fitted.

The reciprocal mean cluster size is computed over disjoint blocks rather than over runs-declustered clusters. This is deliberate. When clusters are taken from runs declustering they partition the exceedances exactly, so the reciprocal mean cluster size reduces algebraically to the runs estimator and carries no independent information; reporting the two as separate estimators would present one quantity twice.

\subsection{Reference extremal indices}

Bias, root mean squared error and coverage all require a reference value of \(\theta\). Closed-form references exist only for the independent systems and for the noiseless logistic map at its periodic target, where \(\theta = 1/2\). Every noisy configuration lies outside those cases, so references there are established numerically from a long realisation of \(4 \times 10^{5}\) samples.

The numerical construction is deliberately conservative, because a reference obtained from one declustering rule is not independent of the estimators being judged: scoring against a runs-based reference would favour the runs estimator by construction. Two structurally different long-run estimates are therefore computed, one by runs declustering at a short run length and one by disjoint blocks at a block scaled to hold roughly one expected exceedance. A numerical reference is accepted only where the two agree within \(0.05\); otherwise the configuration is recorded as one where \(\theta\) could not be established, and it contributes no bias, RMSE or coverage evidence. Both long-run settings are calibrated against the two systems whose value theory already fixes, recovering \(0.495\) against a theoretical \(0.5\) and \(0.961\) against a theoretical \(1.0\).

Theoretical and numerical references are never pooled within a reported row, and every table states which kind it used.

\subsection{Interval estimation}

None of the six estimators has a usable closed-form interval in this setting, so intervals are obtained by a circular block bootstrap over the exceedance indicator series. Resampling whole blocks preserves the short-range dependence that the extremal index measures; an independent bootstrap would destroy exactly the structure under study and return intervals centred on one. Resamples that fail to yield an admissible estimate are counted rather than discarded, so the reported resample success rate is honest.

\subsection{What this grid can and cannot establish}

The simulation evidence is not uniform across the clustering regimes, and the reported scope is limited accordingly. The high-replication grid runs at a single noise level, and at that noise level the numerically established references sit close to independence: of its nine referenced process families, eight have \(\theta \geq 0.9\). Its 500 replications per cell therefore buy precision about \emph{near-independent} processes. Clustered evidence does exist --- the cyclic degradation family is established between \(\theta = 0.049\) and \(\theta = 0.175\), and the smoothing family at \(\theta = 0.272\) --- but the strongly clustered cells carry far fewer replications, so their coverage estimates are correspondingly imprecise.

Accordingly, this study reports the high-replication grid as evidence about near-independent processes, and treats the clustered cells as supporting rather than primary evidence. Estimator failures, non-estimable quantities and unestablished references are preserved in the artifacts as part of that scope statement rather than removed from it.

\section{Datasets}
\label{sec:datasets}

The datasets are assigned roles before interpretation. MetroPT and MetroPT2 are the primary event-level case studies because they contain ordered compressor telemetry and labelled failure periods \citep{Veloso2022MetroPT,MetroPT2Zenodo2023}. They support early-warning, alarm-episode, recurrence, and per-failure questions, but their independent event count is small. SCANIA Component X is an external vehicle-level stress test with irregular fleet histories, repair labels, and censoring \citep{Kharazian2025Scania}. Hydraulic Systems is a supplementary laboratory condition-monitoring demonstration with load-cycle component states \citep{Helwig2018HydraulicSystems}. SECOM is a supplementary high-dimensional production-quality demonstration with wafer-level yield labels \citep{McCann2008SECOM}.

This role assignment prevents a common mistake: row counts do not become independent failures. Table~\ref{tab:dataset-splits} is retained as a descriptive split inventory for the synthetic paper-asset cache. Table~\ref{tab:real-data-characteristics} records real-data source, independence unit, event or target source, and processed row counts. The independent unit column is the primary evidential quantity for empirical interpretation.

\input{../reports/paper/latex/dataset_split_summary.tex}

\begin{table}[t]
  \centering
  \caption{Registered real-data characteristics generated by the real-data matrix.}
  \label{tab:real-data-characteristics}
  \resizebox{\linewidth}{!}{\input{../artifacts/real_data_matrix/dataset_characteristics.tex}}
\end{table}

\begin{table}[t]
  \centering
  \caption{Dataset roles, sampling structures, independent units, verification status, and diagnostic result status.}
  \label{tab:evidence-scope}
  \resizebox{\linewidth}{!}{\input{../artifacts/real_data_matrix/evidence_scope.tex}}
\end{table}

The supported estimands differ. MetroPT-style rows address event-level early warning on temporal telemetry splits. SCANIA addresses vehicle-level repair-risk ranking on a published test split. Hydraulic Systems addresses cycle-level component-condition detection. SECOM addresses production-example yield-failure detection. Hydraulic Systems and SECOM are therefore not interpreted as event-level replications of the compressor task. This distinction is part of CLM-007: broader predictive claims require coherent independent units and comparable targets, not only successful data preparation.

\section{Experimental Protocol}
\label{sec:protocol}

The evaluation protocol is frozen before final experiment execution. It declares the train, validation, calibration, and test partitions; the exclusion and maintenance-boundary fields; warning horizons; alarm merge gap; matching tolerance; matching method; bootstrap unit; primary metrics; secondary metrics; and event utility costs. The protocol hash used by this manuscript is \DynEvtProtocolHash{}.

Primary event metrics are event precision, event recall, event \(F_1\), false-alarm events per operating day, duplicate alarms per failure, median warning lead time, and horizon Brier score. Secondary metrics include pointwise ranking metrics, time under warning, and event utility. Baselines are required to use the same causal feature availability and the same frozen event conversion, so a model is not rewarded for using future information or a more favorable matching rule.

Uncertainty is attached to the independent unit declared by the dataset and protocol. For simulation cells, the unit is the Monte Carlo repetition or registered replicate. For MetroPT-style case studies, the unit is the failure episode or operating interval, not the timestamp. For SCANIA, the relevant unit is the vehicle or repair history defined by the dataset release. Multiplicity is handled by keeping primary metrics and hypotheses predeclared and by routing every interpretive statement through the claim ledger.

\section{Results}
\label{sec:results}

\subsection{Simulation validity and estimator stability}

The first research question asks whether the extremal-index diagnostics are stable enough to interpret. The current simulation artifact is a limited recovery study rather than the full Monte Carlo validation required for a methodological paper. Figure~\ref{fig:simulation-rmse} reports runs-estimator RMSE by threshold quantile and simulation family. The result is useful for finding unstable threshold regions, but it is not sufficient to claim estimator reliability across industrial fault mechanisms.

\begin{figure}[t]
  \centering
  \includegraphics[width=0.78\linewidth]{../reports/paper/figures/simulation_bias_rmse.png}
  \caption{Limited simulation recovery summary. The plotted quantity is runs-estimator RMSE by threshold quantile and simulation family; broader Monte Carlo bias, coverage, and failure-rate validation remains a required extension.}
  \label{fig:simulation-rmse}
\end{figure}

Threshold and run-length sensitivity in Figure~\ref{fig:threshold-stability} and Table~\ref{tab:threshold-stability-summary} show that clustering conclusions depend on \(q\) and \(m\). That supports a negative-result interpretation: regime-conditioned dynamical-EVT diagnostics are sensitive design objects, not model-free evidence of impending failure.

\begin{figure}[t]
  \centering
  \includegraphics[width=0.78\linewidth]{../reports/paper/figures/threshold_stability.png}
  \caption{Threshold and run-length stability diagnostic for the registered recurrence observable.}
  \label{fig:threshold-stability}
\end{figure}

\input{../reports/paper/latex/threshold_stability_summary.tex}

\subsection{Alarm conversion and failed-detection regions}

The second research question asks whether declustering and alarm merging reduce duplicate alarms while preserving event recall. In the registered threshold grid, event recall was zero at all reported quantiles, with false-alarm burdens between 133.92 and 168.48 alarm episodes per operating day. This is a failed-detection region under the current observable, threshold grid, and matching rule; it shows that alarm conversion can destroy utility. The full threshold grid is reported in the supplement rather than as a main-text figure because a zero-recall line plot would add no scientific information.

Across all registered combinations of alarm merge gap and post-onset matching tolerance, event recall, precision, and \(F_1\) remained zero. Relaxing post-onset tolerance did not alter the result, so the full grid is reported in the supplement rather than as a main-text sensitivity figure.

\input{../reports/paper/latex/matching_tolerance_summary.tex}

\subsection{Real-data diagnostic outcomes}

The third and fourth research questions ask whether the registered robust extreme-score diagnostic adds operational value after episode conversion, and why empirical failures occur. Table~\ref{tab:event-level-results} gives the held-out event-level rows for MetroPT and MetroPT2. Each dataset contains one labelled test event in the registered split. Both events are detected, but the detection is operationally weak: MetroPT produces 2527 predicted alarm episodes with precision \(0.00040\), and MetroPT2 produces 849 predicted alarm episodes with precision \(0.00118\). The median lead times, 3346 and 3469 samples respectively, are not persuasive by themselves because they are paired with extreme false-alarm burdens.

Hydraulic Systems, SCANIA Component X, and SECOM are reported separately in Table~\ref{tab:non-event-results} because their independent units are cycles, vehicles, and wafers rather than compressor failure episodes. All three have \(F_1=0\) under the conservative high-quantile diagnostic. These are not hidden failures. They are evidence that a simple extreme-score adapter does not transfer across laboratory cycle states, vehicle-level repair risk, and wafer-level yield labels without changing the estimand and evaluation design. Table~\ref{tab:root-cause} summarizes the dominant explanation supported by each diagnostic row.

\input{../reports/paper/latex/industrial_event_level_results.tex}
\input{../reports/paper/latex/industrial_non_event_results.tex}
\input{../reports/paper/latex/root_cause_summary.tex}

The dominant root-cause classifications are therefore provisional but informative. For MetroPT and MetroPT2, the current evidence points to alarm-conversion failure and insufficient independent failures: recall can be obtained, but only by issuing many alarms around sparse labelled events. For SCANIA, the dominant issue is estimand mismatch: a vehicle-level repair-risk task is not the same as compressor event warning. For Hydraulic Systems and SECOM, the current high-quantile score finds no useful positive-event detection, consistent with non-event labels and signatures that need different feature families or supervised baselines.

\subsection{Score stratification and calibration limits}

Figure~\ref{fig:reliability} and Figure~\ref{fig:split-calibration} are retained as score-stratification diagnostics, not probability calibration evidence. The current model outputs are ranks or anomaly scores. The generated bin tables report mean rank score and observed future-event proportion; the independent positive-event evidence is too limited for a strong calibrated probability claim.

\begin{figure}[t]
  \centering
  \includegraphics[width=0.58\linewidth]{../reports/paper/figures/reliability_diagram.png}
  \caption{Score-stratification diagnostic for ranked horizon scores. Probability interpretation is not claimed.}
  \label{fig:reliability}
\end{figure}

\begin{figure}[t]
  \centering
  \includegraphics[width=0.62\linewidth]{../reports/paper/figures/split_calibration_intervals.png}
  \caption{Split-aware score-stratification intervals for ranked horizon scores.}
  \label{fig:split-calibration}
\end{figure}

\subsection{Metric provenance, baselines, and audit timeline}

Table~\ref{tab:metric-provenance} resolves the apparent conflict between strong baseline validation scores and weak operational event metrics. The baseline scores currently come from the synthetic runner validation split and are pointwise smoke checks. The MetroPT and MetroPT2 values are held-out event-level alarm metrics. They are different quantities and cannot be compared as if they were the same result.

\input{../reports/paper/latex/metric_provenance.tex}
\input{../reports/paper/latex/method_scope_completion.tex}

Figure~\ref{fig:metric-collapse} visualizes the same distinction: pointwise smoke \(F_1\) can be high while event-level precision remains operationally poor after alarm conversion.

\begin{figure}[t]
  \centering
  \includegraphics[width=0.76\linewidth]{../reports/paper/figures/metric_collapse_decomposition.png}
  \caption{Metric provenance diagnostic separating synthetic pointwise validation from real event-level alarm precision.}
  \label{fig:metric-collapse}
\end{figure}

The baseline-family table is therefore moved to the supplement. It documents runner coverage under the shared causal prediction contract, not a completed real-dataset event-level baseline comparison. The event timeline in Figure~\ref{fig:timeline} is retained only as a representative timing-audit display; Table~\ref{tab:timeline-traceability} records that it is not a real MetroPT failure-specific figure.

\begin{figure}[t]
  \centering
  \includegraphics[width=0.95\linewidth]{../reports/paper/figures/event_timeline.png}
  \caption{Representative aligned alarm timeline. The x-axis is elapsed hours from failure onset; the panels show score and threshold, operating regime, alarm episodes after the registered merge rule, warning window, and labelled failure interval. This figure is a timing-audit illustration, not an empirical MetroPT result row.}
  \label{fig:timeline}
\end{figure}

\input{../reports/paper/latex/timeline_traceability.tex}

\section{Discussion}
\label{sec:discussion}

The main value of the workflow is that it makes several failure modes explicit. A low extremal-index estimate may indicate clustered extremes under the chosen threshold and run length, but it can also result from periodicity, persistent operating regimes, preprocessing, missingness, or finite-sample noise. Dynamical EVT supplies the language linking recurrence and extremes, yet the industrial interpretation still requires domain evidence. This paper therefore treats \(\theta\) as a dependence diagnostic and not as a direct mechanical degradation variable.

Event-level evaluation changes the operational reading of anomaly scores. Pointwise score improvements can be unhelpful if they fragment one failure into many duplicate alarms or if they achieve recall only by keeping a machine under warning for a large fraction of exposure. Conversely, an apparently sparse alarm policy can be unacceptable if it fires too late for maintenance action. The generated frontier and timeline make these trade-offs visible, but they do not remove the need for cost, staffing, and repair-window assumptions.

The five-dataset matrix changes the evidential status of the package. The repository is no longer only a synthetic workflow plus one compressor-family case study. It now verifies public raw and processed artifacts for compressor telemetry, truck repair histories, hydraulic test-rig cycles, and semiconductor process examples, and it writes one terminal diagnostic row for each. This broader base is valuable because it forces the workflow to handle different schemas, labels, missingness patterns, entity structures, and independence units.

The negative result in CLM-007 is still important. The current generated assets do not support broad industrial performance claims beyond the evaluated evidence. That statement is not a weakness to hide; it is a boundary condition for honest use of the artifacts. MetroPT and MetroPT2 have limited independent failure evidence; SCANIA changes the target estimand to vehicle-level risk stratification; Hydraulic Systems uses laboratory cycle-level component states; and SECOM uses wafer-level yield labels with missing high-dimensional process measurements. Future empirical work should therefore add additional public or prospective datasets and predeclared baselines before making a broader predictive-maintenance claim.

The reproducibility layer is also a scientific constraint. Every generated figure and table carries provenance, the claim ledger blocks unsupported manuscript injection, and the paper-check command audits citations, claims, generated assets, and compiled PDFs. This structure does not guarantee that every modelling choice is optimal, but it makes unsupported interpretation harder to introduce silently.

\section{Conclusion}
\label{sec:conclusion}

The event-level compressor case studies gave a direct negative answer. MetroPT contained one independent labelled test failure; the registered diagnostic detected it, but only with 2527 predicted alarm episodes, event precision \(0.00040\), and a false-alarm burden of 101.17 alarm episodes per operating day. MetroPT2 also contained one independent labelled test failure; it was detected with 849 predicted alarm episodes, event precision \(0.00118\), and 51.29 false-alarm episodes per operating day. These detections are therefore not operationally persuasive: the event was recovered, but the alarm burden was far beyond what would constitute a useful maintenance signal.

The statistical interpretation is that clustered extremes are dependence diagnostics, not maintenance decisions. Threshold and run-length choices changed the estimated clustering behavior, which means the extremal-index estimate is conditional on the observable, regime definition, threshold, and declustering rule. Low estimated extremal index can indicate recurrence, persistence, smoothing, missingness, periodic control behavior, or regime mixture; it is not direct evidence of degradation unless the target region is known to be failure-proximal and the episode-level alarm policy remains useful after conversion.

The four research questions are answered narrowly. Regime conditioning and threshold grids exposed sensitivity, but they did not establish a stable diagnostic suitable for general deployment. Declustering and alarm merging did not reduce duplicate alarms while preserving event recall in the registered failed-detection grid; all tested merge-gap and post-onset tolerance combinations had zero recall, precision, and \(F_1\). The robust extreme-score diagnostic added limited diagnostic value by revealing alarm-conversion failure and estimand mismatch, but it did not add operational early-warning value. The dominant failure mechanisms are sparse independent failures and alarm burden for MetroPT and MetroPT2, vehicle-level estimand mismatch for SCANIA Component X, cycle-state mismatch for Hydraulic Systems, and yield-target mismatch for SECOM.

Recurrence-based EVT can support maintenance decisions only after stronger conditions are met: an identifiable failure-proximal target region, enough independent failures or entities for uncertainty estimation, leakage-free validation, stable regimes and thresholds, fair comparison against simple baselines, acceptable event recall, acceptable false-alarm burden, useful warning lead time, and prospective or external validation. The present study should therefore be read as evidence about the limits of converting clustered extremes into maintenance alarms, not as evidence of predictive superiority.


\bibliographystyle{plainnat}
\bibliography{references}

\end{document}
